% Short version only: what an implementer needs that the body leaves out
% (paper/reviews/developer.md findings 3, 4 and 8 and its closing list).
% Sources: spec/layertree.md 6.2, 6.5, 9.1 to 9.7, 10.3, 10.5, 10.9,
% 10.11 and 11.2 (0.12.x; the voucher message follows the 0.13.0 rule
% that the source leaf is bound without its amount), spec/channel-leaves.md
% C.4.4, C.5.2, C.5.5 and C.9.3, design/receive.md, msg/ and txn/.
% Sizes at D = 16 dense, FULL, on the model of the evaluation section.
\section{Implementer's Notes}\label{app:impl}

The body says what each operation does and who signs it. This appendix
collects what a wallet, a tower or an operator also has to get right:
how a tree is created and recognised, what each party stores and rolls
back, the Lightning receive as a message sequence, the fee reserve every
unilateral path needs, mempool behaviour, leases and channel leaves.
Section~\ref{sec:ops} serves as the per-operation card (who authorises,
what the witness proves, what the outputs must be); the witness order of
each tapleaf is in spec~8, the builders and signers are in the Go
package \path{txn/} (\texttt{BuildExit}, \texttt{BuildJoin},
\texttt{BuildClaim}, \texttt{SignVoucher}, \texttt{SignExitAuth} and the
rest), and what a co-signer verifies before signing is spec~9.2 step~3
for joins, 9.5 for updates and 9.6 for vouchers.

\subsection{Creating and Recognising a Tree}\label{app:impl-create}

The operator fixes the tree's parameters: the profile, the depth $D$,
the operator key, the dust floor $\mathsf{DUST\_MIN}$, the supported
\tr{JOIN}$_h$ heights and \tr{EXIT}$_K$ sizes, the variants (the
exact-payout and tower exits, \tr{UPDATE}$^{x}$, \tr{CLAIM}$^{x}$,
\tr{REWRITE}$_h$ or \tr{REWRITE}$_h^{x}$, \tr{DETACH}$_h$) and
$\mathsf{tree\_id}$. It compiles the taptree from them and publishes the
parameters with its root $m$ (spec 6.2, 9.1). The recommended
$\mathsf{tree\_id}$ is a tagged hash of the outpoint spent by input~0 of
the creation transaction, which is known before that transaction is
built and can be spent only once, so it is unique per lineage; a chosen
constant is not, and the tree output's own outpoint cannot serve, since
$m$ depends on $\mathsf{tree\_id}$. The creator therefore fixes input~0
before compiling the scripts. A zero-value output is invalid, so the
operator seeds the tree with a leaf of its own, or with initial leaves
whose owners have checked their inclusion before funding (spec 6.5).

Before treating an output as a LayerTree, a member, a voucher holder and
a tower each recompute $m$ from the published parameters and check
$\TreeSpk(r, m)$ against the output for the root $r$ they hold a path
under. The scripts are visible only through $m$, so an operator that
adds one tapleaf of its own (say, ``the operator key spends the
output'') produces an output indistinguishable on-chain from an honest
one; the recomputation, and rejecting any taptree with a leaf outside
the expected set, is the check that custody rests on. A credit-leaf
recipient also confirms that no plain \tr{REWRITE}$_h$ is present, and a
voucher holder that the voucher's $\mathsf{tree\_id}$ is the one the
\tr{CLAIM} scripts carry.

\subsection{State, Backups and Reorganisations}\label{app:impl-state}

Table~\ref{tab:impl-state} lists what each party keeps. Every
transition's witness publishes the paths it used, so a party holding a
path for root $r_k$ derives its path for $r_{k+1}$ from the witness of
that spend alone (spec 9.7); a light client finds the spend with BIP-158
filters, which match the spent output's scriptPubKey. The one piece of
data that is not on-chain is a grafted member's path inside its subtree,
which the operator sends at the join and which a member who loses it can
recover only from the operator or a cohort peer, so it is the backup
that matters (spec 10.9).

\begin{table}[htbp]
\centering
\caption{What each party stores. ``Derives'' is what it updates from each spend of the tree output; on a reorganisation every row is rolled back to the surviving $(\mathsf{txid}, r)$.}
\label{tab:impl-state}
\setlength{\tabcolsep}{3pt}\scriptsize
\newcolumntype{L}[1]{>{\raggedright\arraybackslash}p{#1}}
\begin{tabular}{@{}L{1.5cm}L{4.2cm}L{2.7cm}L{3.0cm}@{}}
\toprule
party & stores & must back up & derives from each spend \\
\midrule
Member & its leaf and path; $r$, $m$, $\mathsf{tree\_id}$ and the published parameters; the txids of its join and of the latest spend & the leaf, and its path inside a grafted subtree & its new path, from the witness \\
Voucher holder & the latest voucher, $\ell_{\mathrm{src}}$ and $\ell_{\mathrm{claim}}$; the source's path; a full-balance voucher or an empty position's path; the source's evictor and lease & the voucher and both leaves & the source's path and $\mathsf{seq}$; whether the lease still outlasts a redemption \\
Tower & the \tr{EXIT}$^{t}$ authorisation, the leaf, its path and $m$; a funded fee coin; for a channel leaf, $(\mathsf{update}_n, \mathsf{settle}_n)$ with their signatures (spec C.9.3) & the authorisation and the latest state & the leaf's path; whether a bound field of the leaf changed \\
Operator & the whole tree: every leaf and path, cohort leases, the vouchers it issued as liabilities & the whole tree & every node; which leases have passed \\
\bottomrule
\end{tabular}
\end{table}

A holder who is not a member has no join-time path to start from: it
follows the tree from creation or asks the operator for the source's
path and checks it against the confirmed root. The operator may be the
party it later falls back against; it can refuse the path but not
falsify it, since a path either folds to the root or does not.

\paragraph{Reorganisations.} Vouchers, paths and roots are keyed by
$(\mathsf{txid}, r)$ and roll back with the chain (spec 10.11). A
transition reorganised out is rebuilt against the surviving root and
rebroadcast; a pending exit is rebuilt the same way after any competing
spend. A \tr{JOIN} reorganised out leaves each funder's input its own.
A reorganised \tr{CLAIM}$^{h}$ puts the swap back at the issuer's
discretion, so the holder rebroadcasts it at once. A member paying
off-chain for a grafted leaf waits for its usual confirmation depth
before paying (spec 9.2).

\paragraph{Messages.} Spec~9 leaves message encoding out of scope, but
the Go implementation carries TLV codecs, on lnd's \texttt{tlv} package, for the
join offer (\path{msg/join.go}), vouchers (\path{msg/voucher.go}),
proofs (\path{msg/proof.go}) and every channel-leaf message
(\path{msg/chan.go}: open, state proposals and signatures, HTLC add,
fulfil and fail, rebind and splice).

\subsection{Receiving over Lightning, Step by Step}\label{app:impl-receive}

The flow of \S\ref{sec:receive-short}, in FULL, with what goes wrong at
each step:
(1)~\emph{Setup.} The operator grafts a credit leaf that it owns, with
the recipient's key as exit key, the operator as evictor and a finite
lease. The recipient checks $m$ (Appendix~\ref{app:impl-create}) and its
path. If the operator never grafts it, nothing has been paid yet.
(2)~\emph{Invoice.} The operator draws the preimage and issues the
invoice, so the payer's receipt proves that the operator was paid. A
recipient that wants the receipt to name it draws preimages and hands
them over in advance, and can then be paid offline only that many times.
(3)~\emph{HTLC.} The payer's HTLC reaches the operator, which settles it
by revealing the preimage. If it fails, no voucher is issued and the
recipient is owed nothing.
(4)~\emph{Voucher.} Only then does the operator sign a cumulative voucher
for the new running total and send it, store and forward if the
recipient is offline. A provider that settles and never issues takes
this one payment; that is the recipient's whole exposure.
(5)~\emph{Top-up.} Before the running total exceeds the credit leaf's
amount, the operator raises the amount by an amount-only
\tr{REWRITE}$_h^{x}$, about $122\vB$ per credit leaf batched per cohort (a changed position
without its $16\vB$ member signature),
or pre-funds the leaf. A top-up voids no voucher, since the message does
not bind the amount (\S\ref{sec:voucher-def}). Pre-funded capital beyond
what has been vouched for is the recipient's on every path the operator
can take alone, and comes back only with the recipient's consent.
(6)~\emph{Retirement.} The operator pays the recipient another way, for
instance with a leaf of its own in the next graft, then advances the
credit leaf's $\mathsf{seq}$ and lowers its amount by one
\tr{UPDATE}$^{x}$ that the recipient co-signs, or one
\tr{REWRITE}$_h^{x}$ for a cohort of credit leaves.
(7)~\emph{Fallback.} If the operator stops cooperating, the recipient
redeems its latest voucher with \tr{CLAIM}$^{x}$ and leaves with the
exact exit, paying both fees from a confirmed coin of its own
(Appendix~\ref{app:impl-fees}).

\subsection{Fee Reserves and Break-Even Balances}\label{app:impl-fees}

Only FULL's fee-from-payout exit can pay its fee out of tree value. A
carve (\tr{CLAIM}, \tr{CLAIM}$^{h}$, \tr{CLAIM}$^{x}$, \tr{SPLIT})
conserves the tree's total; the exact-payout exit, which a credit-leaf
tree, every channel leaf and the MID, MIN and CCV tiers use, pays the
whole amount; and a channel leaf's update and settlement are zero-fee
and are bumped only by an anchor child (spec C.5.5). So a member, a
voucher holder and each party to a channel leaf keep a confirmed fee
coin as a condition of holding the position, and a channel-leaf party's
reserve grows with the HTLC value it has in flight, since a counterparty
can delay the settlement at about $313\vB$ a block. A user with no
on-chain coin at all has no unilateral recourse except, in FULL, the
fee-from-payout exit of an ordinary leaf. Table~\ref{tab:impl-fees}
gives each path's fee at three feerates: a balance smaller than the fee
is not worth recovering alone, and ten times the fee is the balance at
which the path costs 10\% (spec 11.2).

\begin{table}[htbp]
\centering
\caption{The fee of each unilateral path at $D = 16$, FULL, on the transaction model, against the balance it recovers.}
\label{tab:impl-fees}
\setlength{\tabcolsep}{4pt}\scriptsize
\begin{tabular}{@{}lcrrrr@{}}
\toprule
path & fee coin & vB & 2\,sat/vB & 10\,sat/vB & 50\,sat/vB \\
\midrule
\tr{EXIT}$_1$, fee from the payout & no & 921 & 1{,}842 & 9{,}210 & 46{,}050 \\
\tr{EXIT}$_1$ or \tr{EXIT}$_1^{=}$, fee input & yes & 1{,}020 & 2{,}040 & 10{,}200 & 51{,}000 \\
channel-leaf force close & yes & 1{,}681 & 3{,}362 & 16{,}810 & 84{,}050 \\
\tr{CLAIM}$^{x}$ then \tr{EXIT}$_1^{=}$ & yes & 2{,}655 & 5{,}310 & 26{,}550 & 132{,}750 \\
\bottomrule
\end{tabular}
\end{table}

\subsection{Mempool Behaviour}\label{app:impl-mempool}

An exit is bumped by replacement, which a plain exit's owner must
re-sign and an \tr{EXIT}$^{t}$ holder rebuilds with its own fee input,
or by a child spending its change or a pay-to-anchor output, which any
transition may carry at an index past its payouts (spec 4.7). Version-3
(TRUC) transactions are not an option for transitions: TRUC caps an
unconfirmed chain at two, and the chain of spends of the tree output is
the throughput (spec 10.5). When a competing transition confirms first,
the pending one is invalid; the wallet updates its path from the winner's
witness and rebuilds. A graft with several funders cannot be replaced
without every funder signing again, so it carries an operator change
output or an anchor and is bumped by a child (spec 9.2).

\subsection{Leases}\label{app:impl-leases}

A lease is what lets the operator sweep a cohort at $158\vB$ a leaf, and
it has a recurring price. Leases are aligned per cohort, since one later
lease holds up a sweep. A renewal moves a member into a cohort with the
new lease, by an \tr{UPDATE} and a whole-balance \tr{SPLIT}
($2{,}616\vB$), or renews a cohort in place with one \tr{REWRITE}$_6$
at $148\vB$ per member, which needs every renewing member's signature;
a member who does not sign keeps the old lease; a sweep needs every
lease in the cohort to have passed, so the operator evicts it alone or
moves it when its lease passes. So a member pays between $148$ and $2{,}616\vB$ per lease
period it stays, and the operator chooses the period. When a lease
passes, the operator should warn the member and offer renewal (spec
9.4); if it evicts instead, the member sees its whole balance arrive at
its exit key, at the operator's fee.

\subsection{Channel Leaves}\label{app:impl-chan}

A channel leaf is a new channel type, not a BOLT channel: LN-Symmetry
updates signed with \op{TX} and CSFS so that they rebind to whatever
funding outpoint the exit creates. Its funding output does not exist
until the leaf leaves the tree, so there is nothing to announce to the
gossip network; the leaf is reached through its counterparty, normally
the operator. Its tower holds $(\mathsf{update}_n, \mathsf{settle}_n)$,
their signatures and a funded fee coin, and, in a tree with the tower
exit, the leaf's \tr{EXIT}$^{t}$ authorisation, which lets it start a
close as well as finish one (spec C.9.3). The deadline budget, with the
counterparty's stale updates in $\beta_{\mathrm{chan}}$, is in
\S\ref{sec:channels}. A splice is one \tr{UPDATE}$^{x}$ ($1{,}048\vB$)
that replaces the leaf's amount and exit key, with no new funding output
and no period in which two versions of the channel are live (spec
C.5.2); it voids the tower's authorisation, so the owner hands over a
fresh one with the first state of the new epoch.
